% !TEX root = ../main_zh.tex
\chapter{总需求、总供给与供给侧}
\label{ch:adas}

IS--LM 模型把价格水平固定了下来。要补全短期的全貌，我们就得让价格动起来，问一问产出与价格究竟是如何共同决定的。答案是一条由 IS--LM 继承而来的总需求曲线，再加上一条刻画厂商如何对价格作出反应的总供给曲线。二者的交点决定了短期均衡；而价格已经充分调整后的长期均衡，则落在潜在产出处那条垂直的长期供给曲线上。有了这套工具，我们终于可以把一次财政扩张完整地追踪到底——先是乘数效应，再是两次挤出效应——并看清为什么在长期里，需求刺激改变的只是价格水平。

\section{总需求曲线}

把 IS 方程与 LM 方程联立起来，消去其中的利率，剩下的便是价格水平与均衡产出之间的一个单一关系。它背后的机制在上一章末尾已经见过：价格水平上升会使实际货币供给 $M/P$ 缩小，把 LM 曲线向上推，抬高利率，从而压低投资和产出。因此均衡产出是价格水平的一个单调减函数，$Y^{d}=Y(P)$。

\begin{figure}[h]
\centering
\input{figures/fig25}
\caption{总需求曲线的推导。在上图中，价格水平上升降低了实际货币供给，使 LM 曲线上移，产出从 $Y_0$ 减少到 $Y_2$。把每个价格水平与其对应的产出画在一起（下图），便描出了向下倾斜的总需求（AD）曲线。}
\label{fig:macro-25}
\end{figure}

\rmkb[总需求曲线为什么向下倾斜]{总需求曲线\emph{并不是}微观经济学里那个故事——买家看到价格更高就少买。在这里，货币可以是完全中性的、公众可以是完全理性的，曲线却依然向下倾斜：价格水平上升降低了实际货币余额，抬高了利率，从而挤出了投资。这条向下的斜率讲的是货币市场和利率，而不是相对价格。}

\defn{总需求曲线}{
\emph{总需求（aggregate demand，AD）曲线}是所有与 IS--LM 均衡相容的 $(Y,P)$ 点的轨迹。它向下倾斜，是因为价格水平上升会减少实际货币余额，从而抬高利率、压低投资和产出。货币政策与财政政策都会使总需求曲线平移；而价格水平本身的变化，则是沿着曲线移动。}

\section{总供给曲线}

在长期，产出由资本、劳动和技术决定，与价格水平无关，因此长期总供给（LRAS）曲线在潜在产出 $\overline{Y}$ 处垂直。在短期，价格是粘性的，供给会对它作出反应。当价格水平上升快于成本（工资的调整存在滞后）时，利润空间被拉大，厂商便把生产扩张到潜在水平之上，于是短期总供给（SRAS）曲线向上倾斜。它的斜率衡量的是供给对价格有多敏感：经济中闲置的资源越多，SRAS 就越平缓；当资源被充分利用时，SRAS 就几乎垂直。

\begin{figure}[h]
\centering
\begin{tikzpicture}[scale=1.0]
  % colours matching the original: orange LRAS, magenta SRAS, black AD
  \definecolor{lrasclr}{RGB}{240,140,0}
  \definecolor{srasclr}{RGB}{230,0,160}
  % potential output level (where LRAS sits) and equilibrium price level
  \def\Ystar{4.6}
  \def\Pstar{3.0}
  % axes
  \draw[-Latex] (0,0) -- (9.2,0) node[right] {$Y$};
  \draw[-Latex] (0,0) -- (0,6.6) node[above] {$P$};
  % long-run aggregate supply: vertical line at potential output
  \draw[lrasclr,line width=1.2pt] (\Ystar,0.15) -- (\Ystar,6.0)
        node[above] {LRAS};
  % short-run aggregate supply: flat at low Y, steepening (convex, J-shaped),
  % monotone increasing, built to pass through the equilibrium (\Ystar,\Pstar):
  %   p(x) = \Pstar + 0.94*(x-\Ystar) + 0.13*(x-\Ystar)^2  (vertex left of domain)
  \draw[srasclr,line width=1.2pt,smooth,samples=140,domain=1.2:6.4]
        plot (\x,{\Pstar+0.94*(\x-\Ystar)+0.13*(\x-\Ystar)*(\x-\Ystar)});
  \node[srasclr] at (7.05,5.85) {SRAS};
  % aggregate demand: downward-sloping straight line through the equilibrium
  %   line of slope -0.9 through (\Ystar,\Pstar)
  \draw[line width=1.2pt] (1.4,{\Pstar-0.9*(1.4-\Ystar)})
        -- (6.8,{\Pstar-0.9*(6.8-\Ystar)});
  \node at (1.85,5.95) {AD};
  % equilibrium point where all three meet
  \fill (\Ystar,\Pstar) circle (1.7pt);
  % guide line to the price axis
  \draw[densely dashed] (\Ystar,\Pstar) -- (0,\Pstar) node[left] {$P^{*}$};
  % potential-output label on the Y axis
  \node[below=4pt] at (\Ystar,0) {\shortstack{Potential\\ GDP}};
\end{tikzpicture}

\caption{总供给与总需求。长期总供给（LRAS）在潜在产出处垂直；短期总供给（SRAS）向上倾斜；总需求（AD）向下倾斜。在完整的均衡中，三条曲线交于同一点。}
\label{fig:macro-26}
\end{figure}

\state{三线同点的均衡}{
在完整的长期均衡中，长期供给曲线、短期供给曲线和总需求曲线交于同一个点，产出恰好处在潜在水平 $\overline{Y}$。}

\section{财政扩张的完整过程}

现在我们可以把政府购买的增加，同时放进三张图里追踪——凯恩斯交叉、IS--LM 与 AD--AS——看着产出依次经过四个水平。

\begin{figure}[h]
\centering
\input{figures/fig27}
\caption{财政扩张的完整调整过程。黑色：在利率和价格水平都固定的情况下，乘数效应把产出抬到 $Y_1$。红色：考虑进流动性偏好后，利率上升，第一次挤出效应把产出拉回 $Y_2$。蓝色：再让价格水平上升，第二次挤出效应把产出拉回 $Y_3$。在长期，成本上升又把产出送回潜在水平 $Y_0$。整个过程中始终有 $Y_0<Y_3<Y_2<Y_1$。}
\label{fig:macro-27}
\end{figure}

经济最初处于潜在产出 $Y_0$ 的长期均衡。这时政府购买增加。
\begin{enumerate}
\item \emph{乘数效应（$Y_0\to Y_1$）。}计划支出上升，经由乘数作用，产出本会达到 $Y_1$——前提是利率和价格水平都保持固定。（在 IS 曲线确定其水平平移量时，我们固定了利率；在 AD 曲线确定其平移量时，我们固定了价格水平。）
\item \emph{IS--LM，第一次挤出（$Y_1\to Y_2$）。}让利率对上升的货币需求作出反应，$r$ 的上升挤出了私人投资，产出回落到 $Y_2$。
\item \emph{AD--AS，第二次挤出（$Y_2\to Y_3$）。}在 $Y_2$ 处，总需求仍然大于短期供给，于是价格水平上升。更高的价格水平缩小了实际余额，进一步抬高利率，挤出更多投资；产出最终稳定在 $Y_3<Y_2$。这便是\emph{第二次挤出效应}。
\item \emph{从短期到长期（$Y_3\to Y_0$）。}随着时间推移，成本（尤其是工资）向更高的价格水平靠拢，利润空间被压缩，短期供给曲线向左平移。每一次左移都重新造成需求超过供给的局面，又把价格往上推，直到供给与需求在潜在产出处重新相遇。
\end{enumerate}

\state{两次挤出效应}{
第一次挤出：需求扩张抬高货币需求，抬高利率，压低投资。第二次挤出：同一次需求扩张抬高了\emph{价格水平}，价格水平降低实际余额，再次抬高利率，进一步压低投资。乘数因此被削减了两次，对产出的净影响可能很小。}

有几点告诫值得明明白白地说出来。价格\emph{只有}在供给与需求失衡时才会变动。IS--LM 均衡是完整均衡的必要条件，却不是充分条件——价格水平还必须让商品市场上的供给与需求出清。而长期的教训是冷峻的：当供给垂直时，单纯的需求刺激只会抬高价格水平，而不会增加产出。

\rmkb[“长期来看，我们都死了”]{凯恩斯这句俏皮话是一种辩护，而非让步。凯恩斯主义政策瞄准的是\emph{短期}：它是逆周期的，意在把经济迅速拉出衰退的泥潭，而不是在繁荣之上再人为制造一场繁荣。两次挤出效应究竟会不会留下乘数的大部分，取决于各条曲线相对斜率的大小——平缓的 LM 曲线或 AD 曲线，又或者大量闲置的资源，都会让刺激政策的效果基本保留下来。当 LM 曲线平缓到即便产出大幅变动也几乎推不动利率时，经济就处在\emph{流动性陷阱（liquidity trap）}之中，货币政策失灵，而财政政策恰恰最为有效。}

\section{供给侧}

上面的实验全都在移动需求。供给的移动则表现得很不一样。沿着一条给定的短期菲利普斯曲线（Phillips curve），通货膨胀与失业反向变动，因此需求驱动的繁荣是用更高的价格换取更低的失业。而一次负向的\emph{供给}冲击，会打破这种取舍关系。

\begin{figure}[h]
\centering
\begin{tikzpicture}[scale=1.0]
  % axes
  \draw[-Latex] (0,0) -- (8.6,0) node[right] {$Y$};
  \draw[-Latex] (0,0) -- (0,7.6) node[above] {$P$};

  % long-run aggregate supply: vertical line
  \draw[thick] (4.6,0) -- (4.6,7.2) node[above] {$\mathrm{LRAS}$};

  % aggregate demand: downward sloping, fixed
  \draw[thick] (1.4,6.4) -- (7.7,1.4);
  \node[above left=-1pt and -1pt] at (1.55,6.3) {$\mathrm{AD}$};

  % original short-run aggregate supply (upward sloping)
  \draw[thick] (1.7,1.6) -- (7.6,5.5);
  \node[right] at (7.6,5.5) {$\mathrm{AS}$};

  % shifted short-run aggregate supply AS' (parallel, moved left/up)
  \draw[thick] (0.6,3.1) -- (6.5,7.0);
  \node[above] at (6.5,7.0) {$\mathrm{AS}'$};

  % equilibria: original at AS x AD, new at AS' x AD
  % AD: through (1.4,6.4)-(7.7,1.4): slope = (1.4-6.4)/(7.7-1.4) = -0.7937
  %  P = 6.4 - 0.7937*(Y-1.4)
  % AS: through (1.7,1.6)-(7.6,5.5): slope = 0.6610;  P = 1.6 + 0.6610*(Y-1.7)
  % intersection E0: 6.4-0.7937(Y-1.4) = 1.6+0.6610(Y-1.7)
  %  -> Y0 ~ 4.62, P0 ~ 4.13
  % (equilibrium dots and dashed guides intentionally omitted, as in the original)

  % leftward-shift arrow from AS toward AS' (in the blank gap between the two AS lines)
  \draw[-Latex,densely dashed] (6.55,4.85) -- (5.45,5.6);
\end{tikzpicture}

\caption{一次负向供给冲击。短期供给曲线左移：产出下降，价格水平却同时上升——这就是滞胀。}
\label{fig:macro-28}
\end{figure}

当短期供给收缩时，产出下降（失业率上升）与价格水平上升（通货膨胀上升）同时发生——这就是被称为\emph{滞胀（stagflation，即 stagnation 与 inflation 的合成）}的“双输”局面。需求政策对它无能为力：刺激会加剧通胀，收缩则会加剧失业。解药只能来自供给侧——提高厂商生产的意愿和能力。这正是上世纪八十年代供给学派（supply-side school）的逻辑：削减边际税率，恢复人们工作和投资的激励，并利用第~\ref{ch:fiscal}~章给出的那个事实——越过某一点之后，更低的税率未必意味着更少的财政收入。

\begin{figure}[h]
\centering
\begin{tikzpicture}[scale=1.0]
  % axes
  \draw[-Latex] (0,0) -- (8.6,0) node[right] {$Y$};
  \draw[-Latex] (0,0) -- (0,6.6) node[above] {$P$};
  % vertical long-run supply curves (AS vertical: output fixed at full-employment level)
  \def\ASx{3.0}   % original AS
  \def\ASpx{5.4}  % shifted AS'
  \draw[line width=1.1pt] (\ASx,0.0) -- (\ASx,6.1);
  \draw[line width=1.1pt] (\ASpx,0.0) -- (\ASpx,6.1);
  % aggregate demand: downward sloping, crossing both AS curves
  % choose AD so that it intersects AS at (3, 4.0) and AS' at (5.4, 2.2)
  \draw[line width=1.3pt] (1.2,5.5) -- (7.2,1.1);
  % labels for the curves
  \node[above] at (\ASx,6.1) {$\mathrm{AS}$};
  \node[above] at (\ASpx,6.1) {$\mathrm{AS}'$};
  \node[above left] at (1.55,5.25) {$\mathrm{AD}$};
  % rightward shift arrow between the two AS curves
  \draw[-Latex,densely dashed] (\ASx+0.25,5.55) -- (\ASpx-0.25,5.55);
  \node[above] at ({(\ASx+\ASpx)/2},5.6) {growth};
  % (equilibrium dots and dashed guides intentionally omitted, as in the original)
\end{tikzpicture}

\caption{长期供给增长。供给曲线右移：产出上升，价格水平下降——这就是真正经济增长的“双赢”。}
\label{fig:macro-29}
\end{figure}

长期供给的增加——生产能力的真正增长——是上面的镜像：产出上升，价格水平下降。这正是经济增长的深层含义：用更低的价格水平获得更多的产出，是这个模型里唯一的“免费午餐”。

\rmk{分析一次短期供给\emph{冲击}时，我们把长期供给固定不变；短期的图仍然可以把长期的变化作为两张静态截面的比较显示出来。}

\section{短期供给曲线的微观基础}

短期供给为什么会向上倾斜？三个模型给出三种理由，但它们都导向同一个简化式，
\[
Y = \overline{Y} + \alpha\,(P - P^{e}), \qquad \alpha>0,
\]
其中产出只有在价格水平超出预期时，才会高于潜在水平。

\subsection{粘性工资模型}

工人和厂商在价格水平揭晓之前就签下名义工资合同，把名义工资固定为 $W=\omega\, P^{e}$，其中 $\omega$ 是目标实际工资。于是实现的实际工资为
\[
\frac{W}{P} = \omega\,\frac{P^{e}}{P}.
\]
当 $P>P^{e}$ 时，实际工资低于目标，厂商多雇人，产出升到潜在水平之上；只有当 $P=P^{e}$ 时，就业和产出才处在它们的潜在水平。该模型预言实际工资是\emph{逆周期}的——在繁荣中下降——这一点在长期里与数据相违背，不过在较短的时间跨度上，这套机制仍有几分作用。

\subsection{不完全信息模型}

这里所有工资和价格都是完全弹性的，但每个生产者只观察到自己那种商品的名义价格，看不到总体价格水平（要滞后一段时间才能知晓）。供给取决于\emph{相对}价格；由于缺少价格水平的信息，厂商便用自己的预期 $P^{e}$ 来替代。如果总体价格水平上升而厂商还看不到这一点，它就会把自己商品名义价格的上升误当成一个有利的\emph{相对}价格信号，于是多生产，结果总产出超过了潜在水平。价格水平只能事后才知道，这是一个相当符合现实的假设。

\subsection{粘性价格模型}

产出与价格同向变动，但因果既可以从 $P$ 流向 $Y$，也可以同样轻易地从 $Y$ 流向 $P$。每家厂商都想根据价格水平和需求状况来定价，
\[
p = P + \alpha\,(Y - \overline{Y}), \qquad \alpha>0,
\]
即产出高于潜在水平时就多要价。可以随意改价的厂商就照这条规则办。而必须提前承诺价格的厂商（因为菜单成本、长期合同或客户关系），则转而依据预期来定价，
\[
p = P^{e} + \alpha\,(Y^{e} - \overline{Y}),
\]
为简便起见，假设它们预期产出处于潜在水平。设有比例为 $s$ 的厂商价格粘性，其余 $1-s$ 完全弹性。均衡价格水平是两类厂商价格的加权平均，
\[
P = s\,P^{e} + (1-s)\big[P + \alpha(Y-\overline{Y})\big],
\]
整理后得到
\[
P = P^{e} + \left[\frac{(1-s)\,\alpha}{s}\right](Y-\overline{Y}),
\qquad\text{即}\qquad
Y = \overline{Y} + \frac{s}{\alpha(1-s)}\,(P-P^{e}).
\]
如果每家厂商都是弹性的（$s=0$），名义收入的增加就会全部传导到价格上：货币是中性的，供给垂直。正是粘性定价者的存在，才给了供给曲线一个正的斜率。与粘性工资模型不同，粘性价格模型意味着实际工资是\emph{顺周期}的：一次负向需求冲击会让弹性厂商降价、让粘性厂商减产并减少劳动需求，而劳动需求曲线的左移压低了实际工资——这一点与数据吻合得更好。
